\uuid{JoKV}
\titre{Étude de suite}
\auteur{bodin ; hanani ; mzari}
\organisation{exo7}
\licence{CC BY-NC-SA 4.0}
\datecreate{2020-01-20}
\niveau{L1}
\module{Analyse}
\chapitre{Suite}
\sousChapitre{Autre}
\difficulte{3}
\theme{Suites, étude de suite}
\exo7qcmid{lille-1.suites.023}

\contenu{
\begin{question}
\qclassification{121.00}
\qtags{motcle=étude de suite}
On considère les suites de termes généraux $\displaystyle u_n=\sum _{k=1}^n\frac{(-1)^k}{k}$, $\displaystyle v_n=\sum _{k=1}^{2n}\frac{(-1)^k}{k}$ et $\displaystyle w_n=\sum _{k=1}^{2n+1}\frac{(-1)^k}{k}$. Quelles sont les bonnes réponses ?
\begin{answers}
    \good{Les suites $(v_n)$ et $(w_n)$ sont convergentes.}
    \good{La suite $(u_n)$ est convergente.}
    \bad{La suite $(u_n)$ est divergente.}
    \bad{L'une au moins des suites $(v_n)$ ou $(w_n)$ est divergente.}
\end{answers}
\begin{explanations}
On vérifie que $(v_n)$ est décroissante, $(w_n)$ est croissante et que $\displaystyle \lim _{n\to +\infty}(v_n-w_n)=0$. Donc ces deux suites sont adjacentes. En particulier, elles convergent et elles ont la même limite $\ell \in \Rr$. Or $v_n=u_{2n}$ et $w_n=u_{2n+1}$, donc, d'après le théorème des suites extraites, la suite $(u_n)$ converge aussi vers $\displaystyle \ell$.
\end{explanations}
\end{question}
}
